\documentclass[11pt]{article}
\usepackage[T1]{fontenc}
\usepackage{lmodern}
\usepackage{amsmath,amssymb}
\usepackage{booktabs}
\usepackage{array}
\usepackage[a4paper,margin=2.5cm]{geometry}
\usepackage[hidelinks]{hyperref}
\setlength{\parskip}{4pt}
\title{\vspace{-1.6cm}\Large Basics and Definitions\\[2pt]
\large A plain-language companion to the research plan on Lamb waves in a magneto-electro-elastic plate}
\author{Dr.\ Vipin \\ Gurugram University}
\date{30 October 2026}
\begin{document}
\maketitle
\vspace{-6mm}

\section{Three properties, three fields}

A piezoelectric material produces a voltage when it is squeezed, and it changes shape when a voltage is
applied to it. Barium titanate is the standard example and it is used in sensors, in capacitors and in
ultrasonic transducers.

A piezomagnetic material does the same thing with magnetism instead of electricity. Squeeze it and the
magnetisation changes; apply a magnetic field and it changes shape. Cobalt ferrite is the common example.

A magneto-electric material is one in which a magnetic field produces an electric polarisation, or an
electric field produces a magnetisation. Single phases with a useful magnetoelectric effect are rare. What
works is a composite: put the piezoelectric and the piezomagnetic phases together so that the strain passes
from one to the other. The magnetic field strains the ferrite through piezomagnetism, the strain is
transmitted across the interface, and the titanate polarises through piezoelectricity. The chain works
because the strain is the messenger.

The consequence for anyone writing a paper on such a material is important. The magnetoelectric coefficient
of a composite is not a material constant in the way that a density is. It depends on how much of each phase
there is, on the shape of the inclusions, on the interface, and on the temperature. Different ways of
averaging the two phases give different answers, and the spread between them is a real uncertainty, not a
rounding error.

\section{What pyroelectric and pyromagnetic mean}

Heating a piezoelectric crystal changes its polarisation even if nothing mechanical happens, because the
crystal expands and the charges shift. The coefficient that measures this is called the pyroelectric
coefficient. In the same way, heating a piezomagnetic crystal changes its magnetisation, and the coefficient
is called the pyromagnetic coefficient.

Both effects run backwards. An electric field applied across the crystal changes its temperature, and that is
called the electrocaloric effect. A magnetic field changes its temperature, and that is called the
magnetocaloric effect. These backwards effects are the reason a magnetocaloric material can be used in a
refrigerator without any moving gas.

The two directions are not independent. If the material is heated and polarises, then an applied field must
also change its heat content, and the two coefficients are tied by the same thermodynamic reciprocity that
links the Seebeck and Peltier effects in a thermocouple. In an equation of motion, this means that a paper
which keeps the pyroelectric term in the electric displacement but drops the electrocaloric term in the heat
equation is not conserving anything. Our model keeps both, and that is a small point of correctness rather
than a new discovery.

\section{How the two phases are averaged}

If you mix two materials, you need a rule for what the mixture's stiffness, permittivity and permeability
are. The simplest rule is the rule of mixtures, which says that the property is the volume average of the two
phases. It is exact when the two phases are arranged as layers in series or in parallel, and it is the rule
most often used for laminates because a laminate really is layers.

A better rule for a material with inclusions is the Mori and Tanaka estimate, which accounts for the shape of
the inclusions and for the fact that each inclusion feels the average field of its neighbours, not the
applied field. For a stack of two continuous layers, the two rules give the same answer in the layering
direction and different answers in the plane.

In this work both routes are carried. The reason is not indecision. If the thermal effect we are chasing is
smaller than the difference between the two routes, then no experiment could tell the effect from the
uncertainty, and the reader deserves to see that on the same graph rather than in a footnote.

\section{Lamb waves in one paragraph}

A plate is thin compared with its length, so a disturbance inside it bounces between the free faces and
travels along the plate. The guided waves that result come in two families, symmetric, where the faces move
in and out in opposite directions, and antisymmetric, where they move together and the plate bends. Phase
velocity depends on frequency, and the graph of one against the other is the dispersion curve. A real wave
also weakens as it goes, and that weakening is written by giving the wave number an imaginary part: the real
part fixes the wavelength, the imaginary part fixes the decay, and the ratio of the two tells you how lossy
the mode is. Comparing plates of different thickness is done by plotting against frequency times thickness.

\section{The size effect}

Very thin plates do not behave like scaled-down thick plates. Two lengths describe the difference. One
softens, because the stress at a point is taken to depend on the strain in a whole neighbourhood and not
at the point alone. The other stiffens, because a region where the strain changes quickly costs more energy than a
region where it is uniform. Both are written as ratios to the thickness, both are small at ordinary
thicknesses, and both raise the order of the equations, which is why the faces need more conditions.

In a laminated composite the size length is not that of either phase. This is one of the places where the
model is weaker than we would like, and the paper treats the size lengths as a sensitivity rather than as
fitted quantities.

\section{The thermal part, briefly}

Strain changes temperature and temperature changes strain, so a wave in any solid drags a temperature field
and loses some energy to it. In the classical description heat diffuses, and the description allows a
temperature change in one place to be felt everywhere at once in a small way, which is unphysical at high
frequency. The modified theory gives heat a short memory, controlled by a relaxation time, and the practical
effect is that the attenuation stops growing with frequency and levels off.

For the composite there is a further difficulty. The relaxation time of the ceramic and that of the ferrite
are very different, because their thermal conductivities and heat capacities are different. An effective lag
for the stack therefore has to be defined, and the choice has to be defended. This is one of the four risks
listed in the plan.

\section{Glossary and units}

\begin{center}
\small
\begin{tabular}{p{4.2cm}p{8.4cm}p{2.4cm}}
\toprule
term & what it means & unit\\
\midrule
piezoelectricity & strain producing charge and charge producing strain & C\,m$^{-2}$\\
piezomagnetism & strain producing magnetisation and field producing strain & N\,A$^{-1}$m$^{-1}$\\
magnetoelectric effect & magnetic field producing polarisation, or the reverse & s\,m$^{-1}$\\
pyroelectricity & polarisation produced by a change of temperature & C\,m$^{-2}$K$^{-1}$\\
pyromagnetism & magnetisation produced by a change of temperature & N\,A$^{-1}$m$^{-1}$K$^{-1}$\\
electrocaloric, magnetocaloric & temperature change produced by an applied electric or magnetic field & \\
product property & property of the composite that neither phase has on its own & \\
rule of mixtures & averaging the property by volume fraction & \\
Mori and Tanaka & averaging that accounts for inclusion shape and interaction & \\
laminate & stack of layers with the layering along the thickness & \\
mode, symmetric, antisymmetric & the two families of guided waves in a plate & \\
dispersion curve & phase velocity drawn against frequency & \\
wavenumber & phase per metre; complex, the imaginary part is the decay & m$^{-1}$\\
attenuation & decay of amplitude per unit distance & Np\,m$^{-1}$\\
frequency times thickness & comparison coordinate for plates of different thickness & Hz\,m\\
size length & the length at which nonlocal or gradient effects become visible & m\\
phase lag & the short memory given to heat by the modified theory & s\\
\bottomrule
\end{tabular}
\end{center}

\end{document}
